\section*{Hoofdstuk \ref{ch:g10:layered-ocean} --- De gelaagde oceaan}

\begin{solution}{exo:g10:layered-ocean:1}
De \omterm{def:g10:layered-ocean:mixed-layer}{gemengde laag} is de oppervlaktelaag van de oceaan waarin temperatuur en
zoutgehalte nauwelijks met de diepte veranderen; de \omterm{def:g10:layered-ocean:thermocline}{thermocline} is de laag
daaronder, waarin de temperatuur snel daalt met de diepte. Wind en golven mengen
de \omterm{def:g10:layered-ocean:mixed-layer}{gemengde laag} (in de winter geholpen door de afkoeling van het oppervlak).
\end{solution}

\begin{solution}{exo:g10:layered-ocean:2}
\omterm{def:g10:layered-ocean:upwelling}{Opwelling} is het opstijgen van koud water van onder de oppervlaktelaag, ter
vervanging van oppervlaktewater dat door de wind is weggeduwd. Het opstijgende
water brengt minerale voedingsstoffen in de verlichte laag, waar het
fytoplankton ze gebruikt om te groeien; het fytoplankton voedt de rest van de
voedselketen.
\end{solution}

\begin{solution}{exo:g10:layered-ocean:3}
Tropische Grote Oceaan: ongeveer \qty{29}{\degreeCelsius} aan het oppervlak en
\qty{4.5}{\degreeCelsius} op \qty{1000}{m}. Noordelijke IJszee: ongeveer
\qty{-1.4}{\degreeCelsius} aan het oppervlak en dicht bij
\qty{0}{\degreeCelsius} op \qty{1000}{m}.
\end{solution}

\begin{solution}{exo:g10:layered-ocean:4}
Bij \qty{0}{\degreeCelsius} (\qty{1028.1}{kg/m^3} tegen \qty{1023.3}{kg/m^3}
bij \qty{25}{\degreeCelsius}). Zeewater: bij \qty{4}{\degreeCelsius} heeft
zuiver water een dichtheid van \qty{1000.0}{kg/m^3}, zeewater met zoutgehalte
35 ongeveer \qty{1027.8}{kg/m^3} (afgelezen op de grafiek).
\end{solution}

\begin{solution}{exo:g10:layered-ocean:5}
De boei drijft op \qty{1000}{m}, zakt naar \qty{2000}{m} en stijgt dan naar het
oppervlak, terwijl ze onderweg temperatuur, zoutgehalte (uit de geleidbaarheid)
en druk meet; aan het oppervlak verstuurt ze haar profiel via een satelliet en
zakt ze weer. Boeien drijven met de stromingen over de hele oceaan en werken
jarenlang, in alle seizoenen, ver van de scheepvaartroutes.
\end{solution}

\begin{solution}{exo:g10:layered-ocean:6}
$361.84\times10^6 \times 3.682 \approx \qty{1.33e9}{km^3}$, ongeveer
\num{1.33} miljard kubieke kilometer. $100/3682 \approx 0.027$: minder dan
\qty{3}{\percent} van de gemiddelde diepte.
\end{solution}

\begin{solution}{exo:g10:layered-ocean:7}
$22.7 - 0.5 = \qty{22.2}{\degreeCelsius}$, bereikt op \qty{15}{m}: de \omterm{def:g10:layered-ocean:mixed-layer}{gemengde
laag} is ongeveer \qty{15}{m} diep. Tussen \qty{20}{m} en \qty{50}{m}:
$(21.8 - 17.3)/30 = \qty{0.15}{\degreeCelsius}$ per meter.
\end{solution}

\begin{solution}{exo:g10:layered-ocean:8}
De zon verwarmt alleen de oppervlaktelaag, en de opwarming en afkoeling van de
seizoenen verspreiden zich alleen door de \omterm{def:g10:layered-ocean:mixed-layer}{gemengde laag}: onder de \omterm{def:g10:layered-ocean:thermocline}{thermocline} is
het water ervan afgesloten. In de winter koelt het oppervlak af, wordt het
dichter en zakt het in het water eronder, en woelen de stormen de zee om: de
menging reikt diep. In de zomer verwarmt de zon een lichte oppervlaktelaag die op
het koudere water drijft; de zachtere winden kunnen die niet naar beneden duwen,
en de \omterm{def:g10:layered-ocean:mixed-layer}{gemengde laag} blijft dun.
\end{solution}

\begin{solution}{exo:g10:layered-ocean:9}
$0.78/0.18 \approx 4.3$ graden rond \qty{10}{\degreeCelsius};
$0.78/0.06 = 13$ graden rond \qty{0}{\degreeCelsius}. In de koude poolzeeën
tellen kleine veranderingen van het zoutgehalte veel zwaarder voor de dichtheid
dan veranderingen van de temperatuur: daar bepaalt het zoutgehalte de
\omterm{def:g10:layered-ocean:stratification}{gelaagdheid}.
\end{solution}

\begin{solution}{exo:g10:layered-ocean:10}
De temperatuur is dezelfde: het diepere water is dichter omdat het zouter is
(33.3 tegen 28.1). De laag ertussen, waar het zoutgehalte snel verandert met de
diepte, is een \omterm{def:g10:layered-ocean:thermocline}{halocline}; hier is ze ook de \omterm{def:g10:layered-ocean:thermocline}{pycnocline}.
\end{solution}

\begin{solution}{exo:g10:layered-ocean:11}
Naar beneden geduwd behoudt het pakketje zijn temperatuur en zoutgehalte, en dus
zijn dichtheid. Omringd door dichter water is het lichter dan het water dat het
verplaatst: de opwaartse kracht duwt het terug naar zijn eigen niveau. Een
pakketje dat omhoog wordt geduwd, zou dichter zijn dan zijn nieuwe omgeving en
terugzakken. De lagen zijn stabiel.
\end{solution}

\begin{solution}{exo:g10:layered-ocean:12}
Massa: $0.001 \times 1028.1 = \qty{1.0281}{kg}$. Nieuw volume: $1.0281/1023.3
\approx \qty{0.0010047}{m^3}$, ongeveer \qty{1004.7}{mL}: het is ongeveer
\qty{4.7}{mL} uitgezet, \qty{0.47}{\percent}. De laag van \qty{700}{m} zou
$700 \times 0.0047 \approx \qty{3.3}{m}$ dikker worden.
\end{solution}

\begin{solution}{exo:g10:layered-ocean:13}
In de noordelijke Atlantische Oceaan in februari. Daar hebben de bovenste
\qty{150}{m} bijna dezelfde temperatuur, en dus bijna dezelfde dichtheid: ze
mengen kost weinig. In de tropische Grote Oceaan stijgt de dichtheid met $1026.0 - 1021.6 =
\qty{4.4}{kg/m^3}$ tussen \qty{50}{m} en \qty{200}{m}: een sterke \omterm{def:g10:layered-ocean:thermocline}{pycnocline}, en de storm zou licht water naar beneden moeten duwen in veel
dichter water.
\end{solution}

\begin{solution}{exo:g10:layered-ocean:14}
Veel chlorofyl ver van de kusten: de band langs de evenaar in de Grote Oceaan
(en in de Atlantische Oceaan), waar de passaatwinden een \omterm{def:g10:layered-ocean:upwelling}{opwelling} aandrijven; en
de hoge breedten, zoals de noordelijke Atlantische Oceaan, waar de \omterm{def:g10:layered-ocean:mixed-layer}{gemengde laag}
in de winter diep reikt en voedingsstoffen omhoogbrengt. Heel weinig chlorofyl:
de centra van de subtropische gyres, sterk \omterm{def:g10:layered-ocean:stratification}{gelaagd}, waar niets diep water
omhoogbrengt en de verlichte laag zonder voedingsstoffen is geraakt.
\end{solution}

\begin{solution}{exo:g10:layered-ocean:15}
Zakken: $1000/(6 \times 3600) \approx \qty{0.046}{m/s}$, ongeveer
\qty{4.6}{cm/s}. Stijgen: $2000/(10 \times 3600) \approx \qty{0.056}{m/s}$,
ongeveer \qty{5.6}{cm/s}. $36\,000/2 = 18\,000$ metingen, gemiddeld
$2000/18\,000 \approx \qty{0.11}{m}$ uit elkaar.
\end{solution}

\begin{solution}{pb:g10:layered-ocean:1}
\textbf{1.} De gyre, met $24.0 - 19.0 = \qty{5.0}{\degreeCelsius}$.

\textbf{2.} De gyre (\qty{20.5}{\degree}~ZB tegen \qty{13.5}{\degree}~ZB).
Omdat hij verder van de evenaar ligt, zou hij de koelere moeten zijn: de breedten
alleen kunnen niet verklaren waarom hij \qty{5}{\degreeCelsius} warmer is.

\textbf{3.} $24.0 - 0.5 = \qty{23.5}{\degreeCelsius}$, bereikt tussen
\qty{50}{m} (23.6) en \qty{100}{m} (22.6), op ongeveer $50 + 50 \times
0.1/1.0 = \qty{55}{m}$.

\textbf{4.} Nee: de temperatuur daalt gestaag vanaf het oppervlak, met
\qty{1.6}{\degreeCelsius} over de eerste \qty{25}{m} en \qty{5.3}{\degreeCelsius}
over de eerste \qty{100}{m}; de \omterm{def:g10:layered-ocean:thermocline}{thermocline} begint aan het oppervlak.

\textbf{5.} Tussen \qty{200}{m} en \qty{300}{m}: $20.3 - 14.4 =
\qty{5.9}{\degreeCelsius}$ per honderd meter.

\textbf{6.} Voor Peru, tussen \qty{50}{m} en \qty{100}{m}: ongeveer
\qty{70}{m} door interpolatie (het volledige profiel geeft \qty{65}{m}). In de
gyre, tussen \qty{200}{m} en \qty{300}{m}: ongeveer $200 + 100 \times 5.3/5.9
\approx \qty{290}{m}$.

\textbf{7.} Ja, op elk niveau van beide kolommen: lichter water ligt op dichter
water, en beide kolommen zijn stabiel.

\textbf{8.} Zouter water zou dichter moeten zijn; toch is het oppervlaktewater
van de gyre lichter: \qty{1024.71}{kg/m^3} tegen \qty{1025.08}{kg/m^3}.

\textbf{9.} Het water van de gyre is ongeveer 1.3 eenheden zouter, wat ongeveer
$1.3 \times 0.8 \approx \qty{1.0}{kg/m^3}$ toevoegt; het is \qty{5}{\degreeCelsius}
warmer, wat er bij deze temperaturen ongeveer
$5 \times 0.27 \approx \qty{1.4}{kg/m^3}$ af haalt. De temperatuur wint.

\textbf{10.} $1026.44 - 1025.08 = \qty{1.36}{kg/m^3}$.

\textbf{11.} $1025.30 - 1024.71 = \qty{0.59}{kg/m^3}$.

\textbf{12.} Van het oppervlak tot \qty{500}{m} stijgt de dichtheid met
$1026.92 - 1024.71 = \qty{2.21}{kg/m^3}$, waarvan maar \qty{0.59}{kg/m^3}
boven \qty{200}{m}: het grootste deel, \qty{1.62}{kg/m^3}, ligt onder
\qty{200}{m}, in de \omterm{def:g10:layered-ocean:thermocline}{thermocline} tussen \qty{200}{m} en \qty{500}{m}.

\textbf{13.} Voor Peru stijgt de dichtheid met $1026.97 - 1025.08 =
\qty{1.89}{kg/m^3}$ van het oppervlak tot \qty{500}{m}, waarvan
\qty{1.36}{kg/m^3} in de bovenste \qty{200}{m}: het grootste deel zit dicht bij
het oppervlak.

\textbf{14.} Op \qty{500}{m} lijken de twee plaatsen bijna op elkaar
(\qty{7.9}{\degreeCelsius} en \qty{7.1}{\degreeCelsius},
\qty{1026.97}{kg/m^3} en \qty{1026.92}{kg/m^3}): onder allebei ligt hetzelfde
koude water; het verschil is hoe dicht het bij het oppervlak komt.

\textbf{15.} De wind die langs de kust naar de evenaar waait, duwt het
oppervlaktewater naar zee; water van onderen stijgt dicht bij de kust op om zijn
plaats in te nemen. Het koude water van de \omterm{def:g10:layered-ocean:thermocline}{thermocline} wordt omhooggetrokken,
zodat de \omterm{def:g10:layered-ocean:thermocline}{thermocline} naar de kust toe omhoogloopt en bijna tot het oppervlak
reikt.

\textbf{16.} Voor Peru: water van \qty{13}{\degreeCelsius} tot
\qty{15}{\degreeCelsius}, van onder de \omterm{def:g10:layered-ocean:thermocline}{thermocline}, ligt binnen de bovenste
\qty{100}{m}, binnen het bereik van het mengen door de wind en van het licht. In
de gyre ligt het onder \qty{250}{m}, onder een warme, lichte laag die de wind
niet zo diep kan mengen.

\textbf{17.} De \omterm{def:g10:layered-ocean:upwelling}{opwelling} zou stoppen: het oppervlaktewater zou in de zon
opwarmen, er zou zich bovenaan een warme, lichte laag vormen, en de aanvoer van
voedingsstoffen naar de verlichte laag zou slinken.

\textbf{18.} De kust van Peru is lichtblauw tot geel (veel chlorofyl); het
midden van de gyre van de zuidelijke Grote Oceaan hoort bij de donkerblauwste
gebieden van de kaart (heel weinig).

\textbf{19.} Fytoplankton heeft naast licht ook voedingsstoffen nodig. In de gyre
houdt de sterke \omterm{def:g10:layered-ocean:stratification}{gelaagdheid} het diepe water, rijk aan voedingsstoffen, ver onder
de verlichte laag, en zijn de voedingsstoffen van de oppervlaktelaag al lang
opgebruikt.

\textbf{20.} $1.36/0.59 \approx 2.3$: \textbf{de toename van de dichtheid over
de bovenste \qty{200}{m} is voor Peru ongeveer 2.3 keer zo groot als in de
gyre}. De \omterm{def:g10:layered-ocean:thermocline}{pycnocline}, met het koude water vol voedingsstoffen eronder, is voor
Peru tot in de bovenste \qty{200}{m} opgetild, terwijl ze in de gyre dieper ligt,
onder \qty{200}{m}.
\end{solution}
